\documentclass{article}
\usepackage{amsmath}
\usepackage{graphicx}
\graphicspath{ {images/} }
\usepackage[spanish]{babel}
\usepackage{bbm}
\usepackage{amsthm}
\usepackage{authblk}
\usepackage{hyperref}
\usepackage[utf8]{inputenc}
\usepackage{mathrsfs}
\usepackage{amsfonts}
\usepackage{amssymb}
\usepackage{enumerate}
\usepackage[left=3cm,right=2cm,top=2cm,bottom=2cm]{geometry}
\usepackage{epsfig}
 \renewcommand{\baselinestretch}{0.93}
 \thispagestyle{empty}
 \pagestyle{empty}
\setlength{\textheight}{53pc} \setlength{\textwidth} {36pc}
\setlength{\topmargin}{-4mm}
\usepackage{multicol}
\newcommand*{\email}[1]{%
    \normalsize\href{mailto:#1}{#1}\par
    }
\setlength{\columnsep}{1cm}
\newcommand{\N}{\mathbb{N}}
\newcommand{\calM}{\cal{M}}
\newcommand{\calP}{\cal{P}}
\newcommand{\F}{\mathbb{F}}
\newcommand{\R}{\mathbb{R}}
\newcommand{\Z}{\mathbb{Z}}
\newcommand{\Q}{\mathbb{Q}}
\newcommand{\C}{\mathbb{C}}
\newcommand{\bbH}{\mathbb{H}}


\theoremstyle{definition}
\newtheorem{ejer}{Ejercicio}
\author{Marcos Morales}
\affil{Universidad Finis Terrae}
\title{Primera Semana\\}



	


\begin{document}


\begin{picture}(400, 75)
   \put(-15,15){\includegraphics[width=5cm]{UFT.png}}
\end{picture}



\begin{center}
\textbf{\large{Temario Examen}} \\
\end{center}

\setlength{\parindent}{0cm}



\begin{center}
\textbf{Cálculo Vectorial}
\end{center}


- Derivadas Parciales:

\begin{enumerate}
    \item Derivadas parciales básicas
    \item Derivadas de orden superior
\end{enumerate}

- Extremos :
\begin{enumerate}
    \item Puntos Críticos
    \item Máximo local, Mínimo local y punto silla
    \item Matriz Hessiana
    \item Regla de los multiplicadores de Lagrange
\end{enumerate}
 



- Integrales dobles: 

\begin{enumerate}
    \item Integrales sobre regiones generales
\end{enumerate}

- Integrales de línea: 

\begin{enumerate}
    \item Integrales de línea sobre campos escalares y campos vectoriales
    \item Teorema de Green
\end{enumerate}    




\end{document} 